\thispagestyle{empty}
\begin{center}
  \vspace*{1.2em}
  {\small\scshape\color{accent} Technical Report\par}
  \vspace{0.9em}
  {\LARGE\bfseries\color{ink} Trident: Design and Simulation of a\par}
  \vspace{0.3em}
  {\LARGE\bfseries\color{ink} Differential-Drive Pick and Place Robot\par}
  \vspace{1.0em}
  {\color{rule}\rule{0.35\linewidth}{0.6pt}\par}
  \vspace{1.2em}
  {\large Joshua John Leo Valan Johnson\par}
  \vspace{0.3em}
  Roll No.\ 110123044\par
  National Institute of Technology Tiruchirappalli\par
  \vspace{1.0em}
  October 2026\par
  \vspace{1.6em}
  {\large\bfseries\color{ink} Abstract\par}
\end{center}
\vspace{0.2em}

\begin{quote}
\small
Trident is a small differential-drive robot that I designed and built for a competition
task: to follow a line and then to carry coloured coins, each to the bin of its colour. This
report presents the mechanical design of the robot and a simulation of it in which the
whole task is carried out autonomously.

The robot is a printed plate of $180\times150\unit{mm}$ on two N20 gearmotors and a ball
caster, designed in SolidWorks and completed with its electronics in fifteen days. Its
gripper does not lift. Two jaws with a parabolic profile, coupled by a pair of equal
involute gears and driven by one servo, cage a coin on the floor, and the robot slides it to
its destination; ridges on the jaws push the coin clear on release. A sliding condition for
a coin on a parabolic jaw is derived as the design rule for the profile. The mass properties
of the CAD model are then used to check the design. They show that the centre of mass lies
inside the support triangle but that the wheels carry only 18\,\% of the weight, which
limits the acceleration to 1.8\unit{m/s^2} at a friction coefficient of one and leaves a
tipping margin of a factor of two. The depth camera cannot see the last 0.15\unit{m} of
floor in front of the robot, and this constraint shapes the software.

The robot is modelled in Gazebo from the CAD export, with models of the depth camera, the
inertial sensor and an eight-channel infrared array. The floor markings of the mission
arena are dimensioned from the pitch of that array. Mapping, localisation and navigation
use RTAB-Map and Nav2; a line follower, docking sequences and a mission state machine were
written for the robot. The damping of the line follower is shown to be set by the distance
of the array ahead of the axle.

In twelve simulated missions with four obstacle layouts the robot delivered all 36 coins,
in 300\unit{s} per mission on average, scored by a referee on ground truth. The line was
followed within 6\unit{mm}, the pose in the map was accurate to 10\unit{mm} RMS, the robot
stopped on the pick markers within 7\unit{mm}, and coins were placed within 21\unit{mm} of
the centres of their squares.

All measurements were taken in simulation. The printed robot won the district round of the
competition and failed in the state round when its infrared array stopped working; no
measurements exist for the hardware.

\medskip
\noindent\textbf{Keywords:} mobile robot design, differential drive, caging gripper, line
following, visual SLAM, navigation, Gazebo, ROS~2.
\end{quote}
